Лето, письменный стол. Направо дверь. На столе картина. На картине нарисована лошадь, а в зубах у лошади цыган. Ольга Петровна колет дрова. При каждом ударе с носа Ольги Петровны соскакивает пенсне. Евдоким Осипович сидит в креслах и курит.

\textsc{Ольга Петровна} (ударяет колуном по полену, которое, однако, нисколько не раскалывается).

\textsc{Евдоким Осипович}. Тюк!

\textsc{Ольга Петровна}. (Надевая пенсне, бьет по полену).

\textsc{Евдоким Осипович}. Тюк!

\textsc{Ольга Петровна} (Надевая пенсне). Евдоким Осипович! Я вас прошу, не говорите этого слова <<тюк>>.

\textsc{Евдоким Осипович}. Хорошо, хорошо.

\textsc{Ольга Петровна} (Ударяет колуном по полену).

\textsc{Евдоким Осипович}. Тюк!

\textsc{Ольга Петровна} (надевая пенсне). Евдоким Осипович! Вы обещали не говорить этого слова <<тюк>>.

\textsc{Евдоким Осипович}. Хорошо, чорошо, Ольга Петровна! Больше не буду.

\textsc{Ольга Петровна} (Ударяет колуном по полену).

\textsc{Евдоким Осипович}. Тюк!

\textsc{Ольга Петровна} (надевая пенсне) Это безобразие! Взрослый пожилой человек и не понимает простой человеческой просьбы!

\textsc{Евдоким Осипович}. Ольга Петровна! Вы можете спокойно продолжать вашу работу. Я больше мешать не буду.

\textsc{Ольга Петровна}. Ну я прошу вас, я очень прошу вас: дайте мне расколоть хотя бы это полено.

\textsc{Евдоким Осипович}. Колите, конечно, колите!

\textsc{Ольга Петровна} (Ударяет колуном по полену).

\textsc{Евдоким Осипович}. Тюк!

Ольга Петровна роняет колун, открывает рот, но ничего не может сказать. Евдоким Осипович встает с кресел, оглядывает Ольгу Петровну с головы до ног и медленно уходит. Ольга Петровна стоит неподвижно с открытым ртом и смотрит на удаляющегося Евдокима Осиповича.

\centerline{Занавес медленно опускается}

\begin{flushright} <\dots> \end{flushright}
